%% This is csc-poster-sample, a sample CSC poster
%% Copyright (C) 2003  Computer Science Club of the University of Waterloo
%% Copyright (C) 2010  Computer Science Club of Brock University
%%
%% This program is free software; you can redistribute it and/or modify
%% it under the terms of the GNU General Public License as published by
%% the Free Software Foundation; either version 2 of the License, or
%% (at your option) any later version.
%%
%% This program is distributed in the hope that it will be useful,
%% but WITHOUT ANY WARRANTY; without even the implied warranty of
%% MERCHANTABILITY or FITNESS FOR A PARTICULAR PURPOSE.  See the
%% GNU General Public License for more details.
%%
%% You should have received a copy of the GNU General Public License
%% along with this program; if not, write to the Free Software
%% Foundation, Inc., 59 Temple Place, Suite 330, Boston, MA  02111-1307  USA

\documentclass{article}
\usepackage{csc-poster}
\usepackage{url}
%\renewcommand{\familydefault}{\sfdefault} %sffamily

\begin{document}

%% TITLE
% This command lays out the top of the document.
% The first parameter is the title.
% The second parameter is the speaker. This parameter may be empty.
\cschead{ACM-ICPC Tryouts}{}

%% SUBTITLE
% This command prints the sub-title and is completely optional.
%\cscsubtitle{International Collegiate Programming Contest}

We invite you to tryout for a spot on the Brock ACM-ICPC team!\\
\\
The ICPC fosters creativity and teamwork in building new software programs,
and enables students to test their ability to perform under pressure.
It is the oldest, largest, and most prestigious programming contest in the world.\\
\\
% Reconsider including this paragraph
%The contest pits teams of three university students against eight or more
%complex, real-world problems with a grueling five-hour deadline. Huddled around
%a single computer, competitors race against the clock in a battle of logic,
%strategy and mental endurance. Teammates collaborate to rank the difficulty of
%the problems, deduce the requirements, design test cases, and write programs
%that solve the problems under the intense scrutiny of expert judges. The team
%that solves the most problems in the fewest attempts in the least cumulative
%time is declared the winner. 
The annual event is comprised of several levels of competition:\\
\textbf{local contests}, \textbf{regional contests}, and \textbf{world finals}.\\
\\
%\begin{description}
%  \item[Local Contests] \hfill \\
%    Universities choose teams or hold local contests to select one or more teams
%    to represent them at the next level of competition.
%  \item[Regional Contests] \hfill \\ % November 3, 2012, U of Windsor
%    Universities send their teams to compete against other universities in their
%    region. Brock participates in the East Central North America Regionals.
%  \item[World Finals] \hfill \\ % June 30 - July 4, 2013, Saint Petersburg, Russia
%    World finalist teams will compete for awards, prizes and bragging rights.
%\end{description}
%\noindent
Programming language options are \emph{C}, \emph{C++}, and \emph{Java}.\\
Programming experience is strongly recommended but not required.\\
\\
Hardcopy reference materials are allowed (anything on paper).\\
Digital references are \emph{not} allowed (internet, physical calculators, ebooks, etc.)\\
\\
For more information about the ICPC, visit \url{http://cm.baylor.edu}\\
For practice problems, visit \url{http://uva.onlinejudge.org}

%% FOOTER
% The first through third arguments tell us the date and time.
% The last argument is a pithy saying.  It may be empty.
% Following the footer, the CSC logo is printed.
\cscfoot{Local Contest}{Friday, 21 September @ 5:00 - 8:00 PM}{MCD205, Brock University}{Questions? Email \url{csc@cosc.brocku.ca}}

% Make sure to comment out the csclogonowords image in csc-poster.sty, line 144.
\vspace{-15mm}
\begin{center}
\includegraphics[scale=0.6]{icpc_logo.eps}
\end{center}

\end{document}
